% Point to the main file
\documentclass[../../Main.tex]{subfiles}

% ========== Start the document ==========
\begin{document}

\section{Section 7.1 - Intro to Probability}

\subsection{Event and Sample Space}

\begin{definition}[Event and Sample Space]
    Let $S$ be a finite, non-empty set called a \yellowbox{sample space} and a set $E \subseteq S$ called the \yellowbox{event space}.
    The probability of $E$ occuring is defined as

    $$P(E) = \frac{|E|}{|S|}$$
\end{definition}

\begin{example}
    Let's do some examples

    \begin{enumerate}
        \item
        {
            Determine the probability of rolling the following numbers when 2 dice are rolled

            \begin{alphabet}
                \item
                {
                    12

                    \begin{align*}
                        S &= \text{ possible rolls of 2 dice, } |S| = 36 \\
                        E &= \text{ rolls that give a 12, } |E| = 1 \\
                        P(E) &= \frac{|E|}{|S|} \\
                        &= \boxed{\frac{1}{36}} \\
                    \end{align*}
                }
                \item
                {
                    2

                    \begin{align*}
                        S &= \text{ possible rolls of 2 dice, } |S| = 36 \\
                        E &= \text{ rolls that give a 2, } |E| = 1 \\
                        P(E) &= \frac{|E|}{|S|} \\
                        &= \boxed{\frac{1}{36}} \\
                    \end{align*}
                }
                \item
                {
                    3

                    \begin{align*}
                        S &= \text{ possible rolls of 2 dice, } |S| = 36 \\
                        E &= \text{ rolls that give a 3, } |E| = 2, (1 + 2), (2 + 1) \\
                        P(E) &= \frac{|E|}{|S|} \\
                        &= \frac{2}{36} \\
                        &= \boxed{\frac{1}{18}} \\
                    \end{align*}
                }
                \item 
                {
                    4

                    \begin{align*}
                        |S| &= 36 \\
                        E &= \{(1, 3), (2, 2), (3, 1)\} \\
                        |E| &= 3 \\
                        P(E) &= \frac{3}{36} \\
                        &= \boxed{\frac{1}{12}} \\
                    \end{align*}
                }

                This is just

                \begin{tablewithheader}
                    { |c|c| }
                    { Sum of dice & Probability }
                    2  & $\frac{1}{36}$ \\ [1ex]
                    3  & $\frac{2}{36}$ \\ [1ex]
                    4  & $\frac{3}{36}$ \\ [1ex]
                    5  & $\frac{4}{36}$ \\ [1ex]
                    6  & $\frac{5}{36}$ \\ [1ex]
                    7  & $\frac{6}{36}$ \\ [1ex]
                    8  & $\frac{5}{36}$ \\ [1ex]
                    9  & $\frac{4}{36}$ \\ [1ex]
                    10 & $\frac{3}{36}$ \\ [1ex]
                    11 & $\frac{2}{36}$ \\ [1ex]
                    12 & $\frac{1}{36}$ \\ [1ex]
                \end{tablewithheader}
            \end{alphabet}
        }

        \item
        {
            Determine the probability of rolling a 2 when a single die is rolled

            \begin{align*}
                S &= \{1, 2, 3, 4, 5, 6\} \\
                E &= \{2\} \\
                P(\text{rolling 2}) &= \boxed{\frac{1}{6}} \\
            \end{align*}
        }

        \item
        {
            Determine the probability of not rolling a 2 when a die is rolled

            \begin{align*}
                S &= \{1, 2, 3, 4, 5, 6\} \\
                E &= S - \{2\} \\
                &= \{1, 3, 4, 5, 6\} \\
                P(\text{not rolling 2}) &= \boxed{\frac{5}{6}} \\
            \end{align*}

            or

            \begin{align*}
                P(\text{not rolling 2}) + P(\text{rolling 2}) &= 1 \\
                P(\text{not rolling 2}) &= 1 - P(\text{rolling 2}) \\
                &= 1 - \frac{1}{6} \\
                &= \boxed{\frac{5}{6}} \\
            \end{align*}
        } 

        \item
        {
            A card is chosen from a standard 52 deck. Determine the probability that the card is

            \begin{alphabet}
                \item
                {
                    An ace

                    \begin{align*}
                        |S| &= 52 \\
                        |E| &= 4 \\
                        P(\text{ace}) &= \frac{4}{52} \\
                        &= \boxed{\frac{1}{13}} \\
                    \end{align*}
                }

                \item
                {
                    An spade

                    \begin{align*}
                        |S| &= 52 \\
                        |E| &= 13 \\
                        P(\text{spade}) &= \frac{13}{52} \\
                        &= \boxed{\frac{1}{4}} \\
                    \end{align*}
                }

                \item 
                {
                    A spade or an ace

                    \begin{align*}
                        |S| &= 52 \\
                        |E| &= |\text{Ace} \cup \text{Spade}| \\
                        &= |\text{Ace}| + |\text{Spade}| - |\text{Ace} \cap \text{Spade}| \\
                        &= 4 + 13 - 1 \\
                        &= 16 \\
                    \end{align*}

                    \begin{align*}
                        P(\text{ace or spade}) &= \frac{16}{52} \\
                        &= \boxed{\frac{4}{13}} \\
                    \end{align*}
                }
            \end{alphabet}
        }
    \end{enumerate}
\end{example}

\subsection{Probability with Multiple Subsets (OR)}

From these examples we see that if $E_1, E_2 \subseteq S$, then

\begin{align*}
    P(E_1 \cup E_2) &= \frac{|E_1 \cup E_2|}{|S|} \\
    &= \frac{|E_1| + |E_2| - |E_1 \cap E_2|}{|S|} \\
\end{align*}

% ========== End the document ==========
\end{document}